% BEGIN REV09_PELL_MOMENTO_PREVIO_IV
\subsection{Moment orienté de profondeur}
\label{corpus9:iv:pell-momento}

La réalisation du registre de profondeur avec quart de tour
définit le moment paramétrique
\[
 \mathcal C_2(r)=
   \sum_{k\geq0}\frac{(-1)^kr^{2k+1}}{(2k+1)^2},
 \qquad 0<r<1,\qquad
 \mathcal D=r\frac{d}{dr}.
\]
La convergence locale uniforme permet de dériver terme à terme :
\[
 \mathcal D\mathcal C_2(r)=\arctan r,\qquad
 \mathcal D^2\mathcal C_2(r)=\frac{r}{1+r^2}.
\]
En effet, la première série dérivée est l’intégrale à partir de zéro
de la série géométrique \(1/(1+r^2)\). Une seconde dérivation
logarithmique produit la dernière égalité. Les conditions
\(\mathcal C_2(0)=\mathcal D\mathcal C_2(0)=0\) fixent les
constantes d’intégration et donnent
\[
 \mathcal C_2(r)=\int_0^r\frac{\arctan t}{t}\,dt
   =\int_0^r\frac{\log(r/t)}{1+t^2}\,dt.
\]
La domination par \(\sum_{n\geq1}n^{-2}\) conserve la limite
\(\mathcal C_2(1)=G_{\mathrm{Cat}}\). La dérivée seconde
en \(r=1\) est comprise au moyen de sa limite d’Abel.
% END REV09_PELL_MOMENTO_PREVIO_IV
